% GENERATED FILE. Edit canonical lessons, then run npm run generate:books.
% canonical-source-hash: fnv1a64:11862c2d264a8d48
% canonical-lessons: PT-C433-algazarra, PT-C433-ranger, PT-C433-barata, PT-C433-bolor, PT-C433-sujidade

\chapter{Barulho e sujidade}
\label{ch:pt-barulho-e-sujidade}
\hlchaptermodality{433}

\begin{chapteropening}
\textbf{By the end of this chapter:} \emph{I can complain about noise and dirt: a racket, creaking, a cockroach, mould and dirt.}

Complete the last lesson of chapter 433: I can complain about noise and dirt: a racket, creaking, a cockroach, mould and dirt.
\end{chapteropening}

\section[a algazarra]{a algazarra — the racket, the din}

\label{lesson:PT-C433-algazarra}



\begin{quote}
Before the new one: say the Portuguese for to recharge, then the Portuguese for the power cut.
\end{quote}

\subsection*{You'll want to know: a algazarra}
\textbf{a algazarra} — \textquotedblleft{}the racket, the din\textquotedblright{}.

\begin{quote}
Os vizinhos fizeram uma algazarra até às duas da manhã.
\end{quote}

\begin{quote}
\emph{The neighbours made a racket until two in the morning.}
\end{quote}

Often of happy noise: children at play.

\begin{cousinweb}
\emph{algazarra} comes from Arabic with the article \emph{al-}; EtymDB traces it to \emph{ghazāra} 'abundance', but the details are not settled.
\end{cousinweb}

\begin{grammarlens}[title={What you've built}]
One more word for telling, explaining and describing at B1.
\end{grammarlens}

\subsection*{Your turn}
\begin{itemize}
\raggedright
  \item \emph{Say it:} \emph{a algazarra}
  \item \emph{Say it:} \emph{a algazarra}, once more
  \item \emph{Say it:} a sentence of your own with \emph{a algazarra}
  \item \emph{Recall:} say the Portuguese for to recharge, then the Portuguese for the power cut, then say \emph{a algazarra} again
\end{itemize}

\begin{tcolorbox}[breakable,colback=teal!4,colframe=teal!35!black,title={Before you move on}]
What does a algazarra mean? (\textquotedblleft{}The racket, the din\textquotedblright{}.) Say it once more.
\end{tcolorbox}
\section[ranger]{ranger — to creak; to grind (teeth)}

\label{lesson:PT-C433-ranger}



\begin{quote}
Before the new one: say the Portuguese for the power cut, then the Portuguese for the racket.
\end{quote}

\subsection*{You'll want to know: ranger}
\textbf{ranger} — \textquotedblleft{}to creak; to grind (teeth)\textquotedblright{}.

\begin{quote}
A porta do armário range muito.
\end{quote}

\begin{quote}
\emph{The wardrobe door creaks a lot.}
\end{quote}

\emph{Ranger os dentes} — to grind your teeth.

\begin{cousinweb}
\emph{ranger} is usually traced to Latin \emph{ringī} 'to snarl, to show the teeth' (Lewis \& Short), but the origin is not settled.
\end{cousinweb}

\begin{grammarlens}[title={What you've built}]
One more word for telling, explaining and describing at B1.
\end{grammarlens}

\subsection*{Your turn}
\begin{itemize}
\raggedright
  \item \emph{Say it:} \emph{ranger}
  \item \emph{Say it:} \emph{ranger}, once more
  \item \emph{Say it:} a sentence of your own with \emph{ranger}
  \item \emph{Recall:} say the Portuguese for the power cut, then the Portuguese for the racket, then say \emph{ranger} again
\end{itemize}

\begin{tcolorbox}[breakable,colback=teal!4,colframe=teal!35!black,title={Before you move on}]
What does ranger mean? (\textquotedblleft{}To creak; to grind (teeth)\textquotedblright{}.) Say it once more.
\end{tcolorbox}
\section[a barata]{a barata — the cockroach}

\label{lesson:PT-C433-barata}



\begin{quote}
Before the new one: say the Portuguese for the racket, then the Portuguese for to creak.
\end{quote}

\subsection*{You'll want to know: a barata}
\textbf{a barata} — \textquotedblleft{}the cockroach\textquotedblright{}.

\begin{quote}
Encontrámos uma barata na cozinha.
\end{quote}

\begin{quote}
\emph{We found a cockroach in the kitchen.}
\end{quote}

Same spelling as \emph{barata} 'cheap', the feminine of \emph{barato}, but a different word.

\begin{cousinweb}
\emph{barata} is inherited from Latin \emph{blatta} 'cockroach, moth' (EtymDB).
\end{cousinweb}

\begin{grammarlens}[title={What you've built}]
One more word for telling, explaining and describing at B1.
\end{grammarlens}

\subsection*{Your turn}
\begin{itemize}
\raggedright
  \item \emph{Say it:} \emph{a barata}
  \item \emph{Say it:} \emph{a barata}, once more
  \item \emph{Say it:} a sentence of your own with \emph{a barata}
  \item \emph{Recall:} say the Portuguese for the racket, then the Portuguese for to creak, then say \emph{a barata} again
\end{itemize}

\begin{tcolorbox}[breakable,colback=teal!4,colframe=teal!35!black,title={Before you move on}]
What does a barata mean? (\textquotedblleft{}The cockroach\textquotedblright{}.) Say it once more.
\end{tcolorbox}
\section[o bolor]{o bolor — mould}

\label{lesson:PT-C433-bolor}



\begin{quote}
Before the new one: say the Portuguese for to creak, then the Portuguese for the cockroach.
\end{quote}

\subsection*{You'll want to know: o bolor}
\textbf{o bolor} — \textquotedblleft{}mould\textquotedblright{}.

\begin{quote}
Há bolor na parede da casa de banho.
\end{quote}

\begin{quote}
\emph{There is mould on the bathroom wall.}
\end{quote}

\emph{Ganhar bolor} — to go mouldy.

\begin{cousinweb}
The origin of \emph{bolor} is not settled, and EtymDB has no entry for it.
\end{cousinweb}

\begin{grammarlens}[title={What you've built}]
One more word for telling, explaining and describing at B1.
\end{grammarlens}

\subsection*{Your turn}
\begin{itemize}
\raggedright
  \item \emph{Say it:} \emph{o bolor}
  \item \emph{Say it:} \emph{o bolor}, once more
  \item \emph{Say it:} a sentence of your own with \emph{o bolor}
  \item \emph{Recall:} say the Portuguese for to creak, then the Portuguese for the cockroach, then say \emph{o bolor} again
\end{itemize}

\begin{tcolorbox}[breakable,colback=teal!4,colframe=teal!35!black,title={Before you move on}]
What does o bolor mean? (\textquotedblleft{}Mould\textquotedblright{}.) Say it once more.
\end{tcolorbox}
\section[a sujidade]{a sujidade — dirt, dirtiness}

\label{lesson:PT-C433-sujidade}



\begin{quote}
Before the new one: say the Portuguese for the cockroach, then the Portuguese for mould.
\end{quote}

\subsection*{You'll want to know: a sujidade}
\textbf{a sujidade} — \textquotedblleft{}dirt, dirtiness\textquotedblright{}.

\begin{quote}
Havia sujidade debaixo da cama.
\end{quote}

\begin{quote}
\emph{There was dirt under the bed.}
\end{quote}

In Brazil: \emph{a sujeira}.

\begin{cousinweb}
\emph{sujidade} is built on \emph{sujo} 'dirty', which is inherited from Latin \emph{sūcidus} 'juicy, greasy' (EtymDB).
\end{cousinweb}

\begin{grammarlens}[title={What you've built}]
One more word for telling, explaining and describing at B1.
\end{grammarlens}

\subsection*{Your turn}
\begin{itemize}
\raggedright
  \item \emph{Say it:} \emph{a sujidade}
  \item \emph{Say it:} \emph{a sujidade}, once more
  \item \emph{Say it:} a sentence of your own with \emph{a sujidade}
  \item \emph{Recall:} say the Portuguese for the cockroach, then the Portuguese for mould, then say \emph{a sujidade} again
\end{itemize}

\begin{tcolorbox}[breakable,colback=teal!4,colframe=teal!35!black,title={Before you move on}]
What does a sujidade mean? (\textquotedblleft{}Dirt, dirtiness\textquotedblright{}.) Say it once more.
\end{tcolorbox}
